\documentclass[12pt,a4paper]{article}

% ── Paquetes ──────────────────────────────────────────────────────────────────
\usepackage[utf8]{inputenc}
\usepackage[T1]{fontenc}
\usepackage[spanish]{babel}
\usepackage{geometry}
\usepackage{amsmath}
\usepackage{booktabs}
\usepackage{array}
\usepackage{graphicx}
\usepackage{float}
\usepackage{hyperref}
\usepackage{setspace}
\usepackage{parskip}
\usepackage{titlesec}
\usepackage{fancyhdr}
\usepackage{xcolor}
\usepackage{needspace}
\usepackage{etoolbox}

% ── Márgenes ──────────────────────────────────────────────────────────────────
\geometry{top=2.5cm, bottom=2.5cm, left=3cm, right=2.5cm}

% ── Encabezado y pie de página ────────────────────────────────────────────────


\pagestyle{fancy}
\fancyhf{}
\fancyhead[L]{\small Universidad Distrital FJDC}

\fancyhead[R]{\small Muestreo con Reemplazo}
\fancyfoot[C]{\thepage}
\renewcommand{\headrulewidth}{0.4pt}

% ── Interlineado ──────────────────────────────────────────────────────────────
\setstretch{1.15}

% ── Hipervínculos ─────────────────────────────────────────────────────────────
\hypersetup{
  colorlinks=true,
  linkcolor=black,
  urlcolor=blue,
  citecolor=black
}

% ── Penalizaciones de salto de página ─────────────────────────────────────────
\widowpenalty=10000
\clubpenalty=10000
\raggedbottom

% ══════════════════════════════════════════════════════════════════════════════
\begin{document}

% ── PORTADA (1 sola página) ───────────────────────────────────────────────────
\begin{titlepage}
  \centering
  \vspace*{1.5cm}

  {\Large \textbf{Universidad Distrital Francisco José de Caldas}}\\[0.3cm]
  {\large Facultad de Ingeniería}\\[2cm]

  {\LARGE \textbf{Muestreo con Reemplazo:}}\\[0.4cm]
  {\Large Análisis Estadístico Descriptivo}\\[2.5cm]

  \begin{tabular}{rl}
    \textbf{Asignatura:} & Probabilidad y Estadística \\[0.4cm]
    \textbf{Docente:}    & Alberto Acosta Lopez       \\[0.8cm]
    \textbf{Estudiantes:} & Daniel Alejandro Castro Becerra \quad 20242020271 \\
                          & Andres Felipe Casilimas Bazurto \quad 20251020068 \\[0.8cm]
    \textbf{Grupo:}      & 19
  \end{tabular}

  \vfill
  {\large Bogotá D.C., 2025}
\end{titlepage}

% ── TABLA DE CONTENIDOS ───────────────────────────────────────────────────────
\tableofcontents
\newpage

% ══════════════════════════════════════════════════════════════════════════════
\section{Introducción}

El presente informe describe de forma integral el proceso de recolección y análisis estadístico
de una muestra obtenida mediante la técnica de muestreo aleatorio simple con reemplazo,
aplicada sobre una población de $N = 60$ puntajes simulados que representan el desempeño
de estudiantes en una prueba diagnóstica de matemáticas (escala de 0 a 99).

A partir de dicha población se extrae una muestra de $n = 20$ elementos. El análisis incluye la
descripción de la técnica de muestreo, el cálculo de estadísticos descriptivos, la construcción
de tablas de frecuencias y la elaboración de diagramas estadísticos generados digitalmente
con Python. El objetivo es presentar e interpretar los resultados de forma clara, rigurosa y
reproducible, siguiendo las normas de formato APA (7.ª edición).

% ══════════════════════════════════════════════════════════════════════════════
\section{Técnica de Muestreo: Muestreo con Reemplazo}

\subsection{Definición}

El muestreo aleatorio simple con reemplazo (MASR) es un método probabilístico en el
que cada elemento de la población tiene la misma probabilidad de ser seleccionado en cada
extracción, y el elemento es devuelto a la población antes de la siguiente. Como consecuencia
directa, un elemento puede aparecer más de una vez en la muestra final \cite{scheaffer2006}.

La probabilidad de selección en cada extracción se mantiene constante:

\begin{equation}
  P(\text{seleccionar cualquier dato}) = \frac{1}{N} = \frac{1}{60} \approx 0{,}0167 \;\; (1{,}67\,\%)
\end{equation}

Las principales características del MASR son:
\begin{itemize}
  \item Las extracciones son independientes entre sí.
  \item La probabilidad de selección es constante en cada extracción.
  \item Un mismo elemento puede aparecer múltiples veces en la muestra.
\end{itemize}

\subsection{Procedimiento Aplicado}

Dado que la población consta de $N = 60$ datos, se realizaron $n = 20$ extracciones individuales
con reemplazo usando la biblioteca \texttt{random} de Python \cite{python_random}. En cada extracción se asignó un
índice del 1 al 60 a cada dato, se generó un número aleatorio dentro de ese rango, y se
seleccionó el dato correspondiente antes de devolverlo al conjunto.

% ══════════════════════════════════════════════════════════════════════════════
\section{Descripción de la Población}

\subsection{Datos de la Población ($N = 60$)}

Los 60 valores que conforman la población representan puntajes enteros en la escala $[0, 99]$:

\begin{table}[H]
  \centering
  \caption{Población de 60 puntajes en prueba diagnóstica de matemáticas.}
  \begin{tabular}{cccccccccc}
    \toprule
    \multicolumn{10}{c}{\textbf{Puntajes de la Población}} \\
    \midrule
    81 & 14 &  3 & 94 & 35 & 31 & 28 & 17 & 94 & 13 \\
    86 & 94 & 69 & 11 & 75 & 54 &  4 &  3 & 11 & 27 \\
    29 & 64 & 77 &  3 & 71 & 25 & 91 & 83 & 89 & 69 \\
    53 & 28 & 57 & 75 & 35 &  0 & 97 & 20 & 89 & 54 \\
    43 & 35 & 19 & 27 & 97 & 43 & 13 & 11 & 48 & 12 \\
    45 & 44 & 77 & 33 &  5 & 93 & 58 & 68 & 15 & 48 \\
    \bottomrule
  \end{tabular}
\end{table}

\subsection{Estadísticos Descriptivos de la Población}

Para contextualizar los resultados de la muestra, se presentan los estadísticos principales de
la población completa:

\begin{table}[H]
  \centering
  \caption{Estadísticos descriptivos de la población ($N = 60$).}
  \begin{tabular}{lc}
    \toprule
    \textbf{Estadístico} & \textbf{Valor} \\
    \midrule
    Media ($\mu$)             & 46{,}45 \\
    Desviación estándar ($\sigma$) & 30{,}81 \\
    Valor mínimo              & 0       \\
    Valor máximo              & 97      \\
    Rango                     & 97      \\
    \bottomrule
  \end{tabular}
\end{table}

% ══════════════════════════════════════════════════════════════════════════════
\section{Muestra Obtenida}

\subsection{Datos de la Muestra ($n = 20$)}

Aplicando muestreo con reemplazo se obtuvieron los siguientes 20 valores, ordenados de
menor a mayor:

\[
  \{3,\; 3,\; 11,\; 13,\; 15,\; 20,\; 25,\; 28,\; 29,\; 31,\;
    35,\; 35,\; 35,\; 35,\; 57,\; 77,\; 83,\; 91,\; 94,\; 94\}
\]

Los valores 3, 35 y 94 aparecen más de una vez, lo que evidencia la característica fundamental
del muestreo con reemplazo: la posibilidad de repetición.

\subsection{Tabla de Frecuencias}

\begin{table}[H]
  \centering
  \caption{Tabla de frecuencias de la muestra ($n = 20$).}
  \begin{tabular}{ccccc}
    \toprule
    $x_i$ & Frec.\ abs.\ ($f_i$) & Frec.\ rel.\ ($h_i$) & Frec.\ \% & Frec.\ acum. \\
    \midrule
     3  & 2 & 0{,}10 & 10{,}0\,\% &  2 \\
    11  & 1 & 0{,}05 &  5{,}0\,\% &  3 \\
    13  & 1 & 0{,}05 &  5{,}0\,\% &  4 \\
    15  & 1 & 0{,}05 &  5{,}0\,\% &  5 \\
    20  & 1 & 0{,}05 &  5{,}0\,\% &  6 \\
    25  & 1 & 0{,}05 &  5{,}0\,\% &  7 \\
    28  & 1 & 0{,}05 &  5{,}0\,\% &  8 \\
    29  & 1 & 0{,}05 &  5{,}0\,\% &  9 \\
    31  & 1 & 0{,}05 &  5{,}0\,\% & 10 \\
    35  & 4 & 0{,}20 & 20{,}0\,\% & 14 \\
    57  & 1 & 0{,}05 &  5{,}0\,\% & 15 \\
    77  & 1 & 0{,}05 &  5{,}0\,\% & 16 \\
    83  & 1 & 0{,}05 &  5{,}0\,\% & 17 \\
    91  & 1 & 0{,}05 &  5{,}0\,\% & 18 \\
    94  & 2 & 0{,}10 & 10{,}0\,\% & 20 \\
    \midrule
    \textbf{Total} & \textbf{20} & \textbf{1{,}00} & \textbf{100\,\%} & \\
    \bottomrule
  \end{tabular}
  \par\smallskip
  \small\textit{Nota.} $h_i = f_i/n$. El valor 35 concentra el 20\,\% de las observaciones,
  constituyéndose en la moda.
\end{table}

% ══════════════════════════════════════════════════════════════════════════════
\section{Cálculo e Interpretación de Estadísticos Descriptivos}

\subsection{Medidas de Tendencia Central}

\subsubsection{Media Aritmética}

La media es el promedio aritmético de todos los valores:

\begin{equation}
  \bar{x} = \frac{\sum_{i=1}^{n} x_i}{n}
           = \frac{3 + 3 + 11 + \cdots + 94 + 94}{20}
           = \frac{814}{20} = 40{,}70
\end{equation}

\textbf{Interpretación:} En promedio, los estudiantes seleccionados obtuvieron un puntaje de
40,70 puntos sobre 99, lo que corresponde a un nivel de desempeño medio-bajo en la prueba
diagnóstica.

\subsubsection{Mediana}

Como $n = 20$ (par), la mediana es el promedio de los valores en las posiciones 10 y 11 del
conjunto ordenado:

\begin{equation}
  M_e = \frac{x_{10} + x_{11}}{2} = \frac{31 + 35}{2} = 33{,}0
\end{equation}

\textbf{Interpretación:} El 50\,\% de los estudiantes de la muestra obtuvo un puntaje igual o
inferior a 33 puntos.

\subsubsection{Moda}

\begin{equation}
  M_o = 35 \quad (\text{frecuencia} = 4)
\end{equation}

\textbf{Interpretación:} El puntaje 35 es el más frecuente en la muestra. Su alta frecuencia
relativa (20\,\%) puede explicarse como resultado natural del muestreo con reemplazo, que
favoreció la repetición de este valor.

\subsection{Medidas de Dispersión}

\subsubsection{Varianza Muestral}

\begin{equation}
  S^2 = \frac{\sum_{i=1}^{n}(x_i - \bar{x})^2}{n-1}
       = \frac{17\,934{,}20}{19} = 943{,}91
\end{equation}

\subsubsection{Desviación Estándar}

\begin{equation}
  S = \sqrt{S^2} = \sqrt{943{,}91} = 30{,}72
\end{equation}

\textbf{Interpretación:} La desviación estándar de 30,72 puntos sobre una escala de 0 a 99
indica una dispersión considerable. En promedio, los puntajes se alejan 30,72 unidades de la
media, lo que refleja diferencias significativas en el rendimiento académico de los estudiantes
evaluados.

\subsubsection{Rango}

\begin{equation}
  R = x_{\max} - x_{\min} = 94 - 3 = 91
\end{equation}

\textbf{Interpretación:} El recorrido de los datos es amplísimo (91 puntos de 99 posibles),
lo que confirma la alta heterogeneidad de la muestra.

\subsection{Resumen de Estadísticos}

\begin{table}[H]
  \centering
  \caption{Resumen de estadísticos descriptivos de la muestra ($n = 20$).}
  \begin{tabular}{lcc}
    \toprule
    \textbf{Estadístico}   & \textbf{Símbolo} & \textbf{Valor} \\
    \midrule
    Media aritmética       & $\bar{x}$        & 40{,}70  \\
    Mediana                & $M_e$            & 33{,}00  \\
    Moda                   & $M_o$            & 35       \\
    Varianza muestral      & $S^2$            & 943{,}91 \\
    Desviación estándar    & $S$              & 30{,}72  \\
    Rango                  & $R$              & 91       \\
    Valor mínimo           & $x_{\min}$       & 3        \\
    Valor máximo           & $x_{\max}$       & 94       \\
    \bottomrule
  \end{tabular}
\end{table}

% ══════════════════════════════════════════════════════════════════════════════
\section{Diagramas Estadísticos}

Los siguientes gráficos fueron elaborados digitalmente utilizando la biblioteca
\texttt{matplotlib} de Python \cite{hunter2007}.

\subsection{Histograma de la Muestra}

\begin{figure}[htbp]
  \centering
  \includegraphics[width=0.85\textwidth]{histogram.png}
  \caption{Histograma de frecuencias de la muestra ($n = 20$). La línea discontinua roja
  indica la media ($\bar{x} = 40{,}7$) y la línea punto-raya azul indica la mediana
  ($M_e = 33{,}0$).}
\end{figure}

\textbf{Interpretación:} El histograma muestra una distribución con concentración de datos en
los rangos bajos (0--35) y una cola hacia la derecha. La separación entre la media y la mediana
($\bar{x} > M_e$) confirma la \textit{asimetría positiva} (sesgada a la derecha), originada por
valores altos como 77, 83, 91 y 94.

\subsection{Gráfico de Frecuencias y Frecuencia Acumulada}

\begin{figure}[htbp]
  \centering
  \includegraphics[width=0.85\textwidth]{frecuabs.png}
  \caption{Gráfico de barras de frecuencia absoluta y línea de frecuencia acumulada de la
  muestra. El eje izquierdo corresponde a la frecuencia absoluta y el eje derecho a la
  frecuencia acumulada.}
\end{figure}

\textbf{Interpretación:} La barra más alta corresponde al valor 35 (frecuencia $= 4$),
confirmando la moda. La curva acumulada muestra que el 70\,\% de las observaciones se
concentran en valores iguales o inferiores a 35.

\subsection{Distribución de la Población con Muestra Destacada}

\begin{figure}[htbp]
  \centering
  \includegraphics[width=0.85\textwidth]{poblacion.png}
  \caption{Diagrama de puntos de la población ($N = 60$). Los rombos naranjas indican los
  elementos seleccionados en la muestra. La línea horizontal representa la media poblacional
  ($\mu = 46{,}45$).}
\end{figure}

\textbf{Interpretación:} La muestra con reemplazo cubrió tanto valores bajos como altos,
aunque con tendencia a seleccionar puntajes por debajo de la media poblacional
(40,70 vs.\ 46,45), diferencia que puede atribuirse al azar propio del proceso de muestreo.

% ══════════════════════════════════════════════════════════════════════════════
\section{Conclusiones}

\begin{enumerate}
  \item La técnica de muestreo aleatorio simple con reemplazo garantizó la independencia de
        cada extracción y la constancia de la probabilidad de selección
        ($1/60 \approx 1{,}67\,\%$), tal como lo establece la teoría \cite{scheaffer2006}.

  \item La presencia de valores repetidos (3, 35 y 94) es una consecuencia natural y esperada
        del MASR. Esta característica distingue esta técnica del muestreo sin reemplazo.

  \item La media muestral ($\bar{x} = 40{,}70$) es superior a la mediana ($M_e = 33{,}0$),
        lo que indica una distribución con asimetría positiva: existe un grupo reducido de
        puntajes altos que desplaza el promedio hacia arriba. Esto es consistente con el
        histograma (Figura 1).

  \item La desviación estándar ($S = 30{,}72$) refleja una dispersión alta en relación con la
        escala de medición, señalando diferencias significativas en el nivel de conocimientos
        matemáticos de los estudiantes al inicio del semestre.

  \item La media muestral (40,70) subestima ligeramente la media poblacional (46,45),
        ilustrando el concepto de \textit{error de muestreo}: las estimaciones basadas en
        muestras pueden diferir del parámetro real debido al azar.
\end{enumerate}

% ══════════════════════════════════════════════════════════════════════════════
\section*{Referencias}
\addcontentsline{toc}{section}{Referencias}

\begin{thebibliography}{9}

\bibitem{scheaffer2006}
  Scheaffer, R.~L., Mendenhall, W., \& Ott, L. (2006).
  \textit{Elementos de muestreo} (6.ª ed.).
  Thomson.

\bibitem{walpole2012}
  Walpole, R.~E., Myers, R.~H., \& Myers, S.~L. (2012).
  \textit{Probabilidad y estadística para ingeniería y ciencias} (9.ª ed.).
  Pearson.

\bibitem{hunter2007}
  Hunter, J.~D. (2007).
  Matplotlib: A 2D graphics environment.
  \textit{Computing in Science \& Engineering}, \textit{9}(3), 90--95.
  \url{https://doi.org/10.1109/MCSE.2007.55}

\bibitem{python_random}
  Python Software Foundation. (s.f.).
  \textit{random --- Generate pseudo-random numbers}.
  Python Documentation.
  \url{https://docs.python.org/3/library/random.html}

\end{thebibliography}

\end{document}